\subsection{作用的正向极限}

假设给定相对于同一指标集 I 的两个集合正向系统 \((\Omega_{\alpha}, \phi_{\beta\alpha})\) 与 \((\mathrm{E}_{\alpha}, f_{\beta\alpha})\) ，并且对每个 \(\alpha \in I\) 给定 \(\Omega_{\alpha}\) 在 \(E_{\alpha}\) 上的一个作用。假设

\[
f _ {\beta \alpha} (\omega_ {\alpha}. x _ {\alpha}) = \phi_ {\beta \alpha} (\omega_ {\alpha}). f _ {\beta \alpha} (x _ {\alpha})
\]

对 \(\alpha \leqslant \beta\) 、 \(\omega_{\alpha} \in \Omega_{\alpha}\) 与 \(x_{\alpha} \in \mathrm{E}_{\alpha}\) 。那么所考虑的这族作用称为一个作用正向系统。如同命题2那样可以验证，存在作用 \(h\) —— \(\Omega = \varinjlim \Omega_{\alpha}\) 在 \(\mathrm{E} = \varinjlim \mathrm{E}_{\alpha}\) 上的作用——其描述如下：设 \(\omega \in \Omega\) 与 \(x \in \mathrm{E}\) ；选取 \(\alpha \in \mathrm{I}\) 及 \(\omega_{\alpha} \in \Omega_{\alpha}\) , \(x_{\alpha} \in \mathrm{E}_{\alpha}\) ，使得 \(\omega = \phi_{\alpha}(\omega_{\alpha})\) 与 \(x = f_{\alpha}(x_{\alpha})\) （ \(\phi_{\alpha}: \Omega_{\alpha} \to \Omega\) 与 \(f_{\alpha}: \mathrm{E}_{\alpha} \to \mathrm{E}\) 表示典范映射）；则 \(\omega . x = f_{\alpha}(\omega_{\alpha}. x_{\alpha})\) 。 \(\Omega\) 在 \(\mathrm{E}\) 上的这个作用称为诸 \(\Omega_{\alpha}\) 在诸 \(\mathrm{E}_{\alpha}\) 上作用的正向极限。

若诸 \(\Omega_{\alpha}\) 是幺半群，且每个 \(\Omega_{\alpha}\) 在 \(E_{\alpha}\) 上的作用是一个运算，则正向极限作用是一个运算。

留给读者定义带算子的群的正向系统的正向极限，并验证该极限是带算子的群。

